% Algebra 2 Lesson 2-1 fast-track packet: Vertex Form of a Quadratic Function
% (enVision A2 2024 TE, Example 2 and Example 4 only; Examples 1, 3, 5 cut).
% Proposal for the PLC (docs/a2-fast-lesson-spec.md): Example 2 teaches quiz item 2 (bird's greatest
% height = vertex), Example 4 teaches quiz items 3 and 4 (vertex + one point -> a, then y = ax^2+bx+c).
% Quiz numbers are the Savvas item numbers; the short form keeps them (items 2, 3, 4; 1 and 5 left off).
% DOK 3 kept: Practice & Problem Solving 15 (Amaya's serve, Example 2), as the Exit block.
% No Practice block (R5): Try It 2 rehearses item 2, Try It 4 rehearses items 3 and 4.
% Shared preamble for the "fast" (hacked-down) Algebra 2 lesson packets, Oct 2026.
% Each packet: % Shared preamble for the "fast" (hacked-down) Algebra 2 lesson packets, Oct 2026.
% Each packet: % Shared preamble for the "fast" (hacked-down) Algebra 2 lesson packets, Oct 2026.
% Each packet: \input{fast-packet-template} then \begin{document} ... \end{document}.
\documentclass[12pt]{article}
\usepackage[margin=0.7in]{geometry}
\usepackage{amsmath,amssymb}
\usepackage{booktabs,tabularx,array}
\usepackage{enumitem}
\usepackage[T1]{fontenc}
\usepackage[utf8]{inputenc}
\usepackage{lmodern}
\usepackage{tikz}
\usepackage{pgfplots}
\pgfplotsset{compat=1.18}
\usepackage{xcolor}
\setlist{nosep,leftmargin=1.4em}
\setlength{\parindent}{0pt}
\setlength{\parskip}{6pt}
\renewcommand{\arraystretch}{1.35}
\newcolumntype{Y}{>{\raggedright\arraybackslash}X}
\pagestyle{empty}

% blanks and workspace
\newcommand{\blank}[1][2.6cm]{\rule[-2pt]{#1}{0.5pt}}
\newcommand{\lines}[1][1]{\par\vspace{2pt}\foreach \i in {1,...,#1}{\rule{0pt}{22pt}\rule{\linewidth}{0.4pt}\par}}
\newcommand{\wspart}[2]{\vspace{4pt}{\large\bfseries #1: #2}\par\vspace{-6pt}\rule{\textwidth}{0.4pt}\par}

% header: \packethead{Lesson 1-3}{Piecewise-Defined Functions}
\newcommand{\packethead}[2]{%
{\LARGE\bfseries #2}\hfill Name \blank[5cm]\\[2pt]
Algebra 2 \quad #1 (fast track) \hfill Date \blank[2.5cm] \quad Period \blank[1.2cm]\par\vspace{2pt}}

% boxed goals: \goals{learning goal}{language goal}
\newcommand{\goals}[2]{%
\fbox{\parbox{\dimexpr\textwidth-2\fboxsep-2\fboxrule}{%
\textbf{Learning goal:} #1\\[2pt]
\textbf{Language goal:} #2}}\par}

% quiz target box: \quiztarget{text}
\newcommand{\quiztarget}[1]{%
\vspace{4pt}\colorbox{gray!12}{\parbox{\dimexpr\textwidth-2\fboxsep}{\textbf{On the quiz:} #1}}\par}

% IXL table row helper: code & skill & which quiz question it feeds
\newcommand{\ixlrow}[3]{\texttt{#1} & #2 & #3 \\}

% one grid style for every graph; override xmin/xmax/... per plot
\pgfplotsset{sheetgrid/.style={
  axis lines=middle, grid=both, grid style={gray!45},
  xmin=-10.5, xmax=10.5, ymin=-10.5, ymax=10.5,
  xtick={-10,-8,...,10}, ytick={-10,-8,...,10},
  x=0.28cm, y=0.28cm,
  ticklabel style={font=\scriptsize, fill=white, inner sep=1pt},
  xlabel={$x$}, ylabel={$y$},
  axis line style={thick}, clip=true}}
% unit grid for small windows
\pgfplotsset{unitgrid/.style={sheetgrid,
  xmin=-5.5, xmax=5.5, ymin=-5.5, ymax=5.5,
  xtick={-5,...,5}, ytick={-5,...,5}, x=0.5cm, y=0.5cm}}
% open and closed endpoints
\tikzset{closed/.style={circle, fill, inner sep=1.6pt},
         open/.style={circle, draw, fill=white, inner sep=1.4pt}}
 then \begin{document} ... \end{document}.
\documentclass[12pt]{article}
\usepackage[margin=0.7in]{geometry}
\usepackage{amsmath,amssymb}
\usepackage{booktabs,tabularx,array}
\usepackage{enumitem}
\usepackage[T1]{fontenc}
\usepackage[utf8]{inputenc}
\usepackage{lmodern}
\usepackage{tikz}
\usepackage{pgfplots}
\pgfplotsset{compat=1.18}
\usepackage{xcolor}
\setlist{nosep,leftmargin=1.4em}
\setlength{\parindent}{0pt}
\setlength{\parskip}{6pt}
\renewcommand{\arraystretch}{1.35}
\newcolumntype{Y}{>{\raggedright\arraybackslash}X}
\pagestyle{empty}

% blanks and workspace
\newcommand{\blank}[1][2.6cm]{\rule[-2pt]{#1}{0.5pt}}
\newcommand{\lines}[1][1]{\par\vspace{2pt}\foreach \i in {1,...,#1}{\rule{0pt}{22pt}\rule{\linewidth}{0.4pt}\par}}
\newcommand{\wspart}[2]{\vspace{4pt}{\large\bfseries #1: #2}\par\vspace{-6pt}\rule{\textwidth}{0.4pt}\par}

% header: \packethead{Lesson 1-3}{Piecewise-Defined Functions}
\newcommand{\packethead}[2]{%
{\LARGE\bfseries #2}\hfill Name \blank[5cm]\\[2pt]
Algebra 2 \quad #1 (fast track) \hfill Date \blank[2.5cm] \quad Period \blank[1.2cm]\par\vspace{2pt}}

% boxed goals: \goals{learning goal}{language goal}
\newcommand{\goals}[2]{%
\fbox{\parbox{\dimexpr\textwidth-2\fboxsep-2\fboxrule}{%
\textbf{Learning goal:} #1\\[2pt]
\textbf{Language goal:} #2}}\par}

% quiz target box: \quiztarget{text}
\newcommand{\quiztarget}[1]{%
\vspace{4pt}\colorbox{gray!12}{\parbox{\dimexpr\textwidth-2\fboxsep}{\textbf{On the quiz:} #1}}\par}

% IXL table row helper: code & skill & which quiz question it feeds
\newcommand{\ixlrow}[3]{\texttt{#1} & #2 & #3 \\}

% one grid style for every graph; override xmin/xmax/... per plot
\pgfplotsset{sheetgrid/.style={
  axis lines=middle, grid=both, grid style={gray!45},
  xmin=-10.5, xmax=10.5, ymin=-10.5, ymax=10.5,
  xtick={-10,-8,...,10}, ytick={-10,-8,...,10},
  x=0.28cm, y=0.28cm,
  ticklabel style={font=\scriptsize, fill=white, inner sep=1pt},
  xlabel={$x$}, ylabel={$y$},
  axis line style={thick}, clip=true}}
% unit grid for small windows
\pgfplotsset{unitgrid/.style={sheetgrid,
  xmin=-5.5, xmax=5.5, ymin=-5.5, ymax=5.5,
  xtick={-5,...,5}, ytick={-5,...,5}, x=0.5cm, y=0.5cm}}
% open and closed endpoints
\tikzset{closed/.style={circle, fill, inner sep=1.6pt},
         open/.style={circle, draw, fill=white, inner sep=1.4pt}}
 then \begin{document} ... \end{document}.
\documentclass[12pt]{article}
\usepackage[margin=0.7in]{geometry}
\usepackage{amsmath,amssymb}
\usepackage{booktabs,tabularx,array}
\usepackage{enumitem}
\usepackage[T1]{fontenc}
\usepackage[utf8]{inputenc}
\usepackage{lmodern}
\usepackage{tikz}
\usepackage{pgfplots}
\pgfplotsset{compat=1.18}
\usepackage{xcolor}
\setlist{nosep,leftmargin=1.4em}
\setlength{\parindent}{0pt}
\setlength{\parskip}{6pt}
\renewcommand{\arraystretch}{1.35}
\newcolumntype{Y}{>{\raggedright\arraybackslash}X}
\pagestyle{empty}

% blanks and workspace
\newcommand{\blank}[1][2.6cm]{\rule[-2pt]{#1}{0.5pt}}
\newcommand{\lines}[1][1]{\par\vspace{2pt}\foreach \i in {1,...,#1}{\rule{0pt}{22pt}\rule{\linewidth}{0.4pt}\par}}
\newcommand{\wspart}[2]{\vspace{4pt}{\large\bfseries #1: #2}\par\vspace{-6pt}\rule{\textwidth}{0.4pt}\par}

% header: \packethead{Lesson 1-3}{Piecewise-Defined Functions}
\newcommand{\packethead}[2]{%
{\LARGE\bfseries #2}\hfill Name \blank[5cm]\\[2pt]
Algebra 2 \quad #1 (fast track) \hfill Date \blank[2.5cm] \quad Period \blank[1.2cm]\par\vspace{2pt}}

% boxed goals: \goals{learning goal}{language goal}
\newcommand{\goals}[2]{%
\fbox{\parbox{\dimexpr\textwidth-2\fboxsep-2\fboxrule}{%
\textbf{Learning goal:} #1\\[2pt]
\textbf{Language goal:} #2}}\par}

% quiz target box: \quiztarget{text}
\newcommand{\quiztarget}[1]{%
\vspace{4pt}\colorbox{gray!12}{\parbox{\dimexpr\textwidth-2\fboxsep}{\textbf{On the quiz:} #1}}\par}

% IXL table row helper: code & skill & which quiz question it feeds
\newcommand{\ixlrow}[3]{\texttt{#1} & #2 & #3 \\}

% one grid style for every graph; override xmin/xmax/... per plot
\pgfplotsset{sheetgrid/.style={
  axis lines=middle, grid=both, grid style={gray!45},
  xmin=-10.5, xmax=10.5, ymin=-10.5, ymax=10.5,
  xtick={-10,-8,...,10}, ytick={-10,-8,...,10},
  x=0.28cm, y=0.28cm,
  ticklabel style={font=\scriptsize, fill=white, inner sep=1pt},
  xlabel={$x$}, ylabel={$y$},
  axis line style={thick}, clip=true}}
% unit grid for small windows
\pgfplotsset{unitgrid/.style={sheetgrid,
  xmin=-5.5, xmax=5.5, ymin=-5.5, ymax=5.5,
  xtick={-5,...,5}, ytick={-5,...,5}, x=0.5cm, y=0.5cm}}
% open and closed endpoints
\tikzset{closed/.style={circle, fill, inner sep=1.6pt},
         open/.style={circle, draw, fill=white, inner sep=1.4pt}}

\begin{document}

\packethead{Lesson 2-1}{Vertex Form of a Quadratic Function}
\goals{I can find the vertex, axis of symmetry, maximum or minimum, domain and range of a quadratic function in vertex form, and write its equation from the vertex and one other point.}{I can explain how vertex form describes a parabola using the words vertex, axis of symmetry, maximum, minimum and opens upward or downward.}

\quiztarget{You will (1) tell how $g(x)=(x+2)^2+4$ moves $y=x^2$ (full Savvas quiz only), (2) find a parabola's highest point from vertex form (Example 2), (3) write vertex form from the vertex and a point (Example 4), (4) read the vertex and $y$-intercept off a graph, find $a$, write $y=ax^2+bx+c$ (Example 4), (5) write a shifted parabola as $y=ax^2+bx+c$ (full Savvas quiz only).}

\wspart{Example 2}{Determine Key Features of a Quadratic Function}
\textbf{Vertex form:} $y=a(x-h)^2+k$, $a\neq0$. \ The \textbf{vertex} is $(h,k)$; the \textbf{axis of symmetry} is the line $x=h$. If $a>0$ the parabola \textbf{opens upward} and $k$ is the \textbf{minimum}; if $a<0$ it \textbf{opens downward} and $k$ is the \textbf{maximum}. The domain is always all real numbers.

\begin{minipage}[t]{0.64\textwidth}
\vspace{0pt}
What are the key features of the quadratic function\\ $f(x)=2(x-3)^2+4$?
\begin{enumerate}[label=\textbf{\arabic*.}, itemsep=3pt]
\item Match it to $a(x-h)^2+k$:\quad $a=2$, \ $h=3$, \ $k=4$.
\item \textbf{Vertex:} $(h,k)=(3,4)$. \quad \textbf{Axis of symmetry:} $x=3$.
\item The 2 indicates that the graph \textbf{opens upward} ($2>0$) and is vertically stretched ($|2|>1$).
\item Opens upward, so the vertex is the lowest point: \textbf{minimum} $4$.
\item \textbf{Range:} $\{y\mid y\geq 4\}$. There are no restrictions on the value of $x$, so the \textbf{domain} is all real numbers.
\end{enumerate}
\end{minipage}\hfill
\begin{minipage}[t]{0.33\textwidth}
\vspace{0pt}
\begin{center}
\begin{tikzpicture}
\begin{axis}[sheetgrid, xmin=-0.5, xmax=6.5, ymin=-0.5, ymax=10.5,
  xtick={0,2,4,6}, ytick={0,2,...,10}, minor tick num=1, minor grid style={gray!45},
  x=0.62cm, y=0.36cm]
\addplot[very thick, red, domain=1.27:4.73, samples=80, <->] {2*(x-3)^2+4};
\addplot[thick, red, dashed, <->] coordinates {(3,-0.4) (3,10.4)};
\node[closed, red] at (axis cs:3,4) {};
\node[font=\scriptsize, inner sep=1pt, anchor=north east] at (axis cs:0,0) {0};
\node[font=\scriptsize, fill=white, inner sep=1pt, anchor=west] at (axis cs:3.3,3.2) {$(3,4)$};
\node[font=\scriptsize, fill=white, inner sep=1pt, anchor=west] at (axis cs:3.2,1.4) {$x=3$};
\node[font=\scriptsize, red, fill=white, inner sep=1pt] at (axis cs:5.4,9.2) {$f$};
\end{axis}
\end{tikzpicture}
\end{center}
\end{minipage}

\textbf{Say it:} The vertex of $f$ is \blank[1.6cm], so the axis of symmetry is \blank[1.6cm]. The parabola opens \blank[2cm] because $a=\blank[1cm]$ is \blank[2cm], so $4$ is the \blank[2.2cm] value of $f$.

\wspart{Try It 2}{Key features from vertex form}
Identify the vertex, axis of symmetry, minimum or maximum, domain, and range of the function $f(x)=-(x+4)^2-5$.

Hint: $x+4=x-(\blank[0.8cm])$, so $h=\blank[0.8cm]$.

\renewcommand{\arraystretch}{1.8}
\begin{tabularx}{\textwidth}{|l|Y|l|Y|}
\hline
$a=$ & & opens upward or downward? & \\ \hline
vertex $(h,k)$ & & axis of symmetry & \\ \hline
minimum or maximum? & & value & \\ \hline
domain & & range & \\ \hline
\end{tabularx}
\renewcommand{\arraystretch}{1.35}

\newpage

\wspart{Example 4}{Write an Equation of a Parabola Given the Graph}
The height of a hit ball is a quadratic function of the time it has been in the air. The ball is hit with an initial height of 4 feet and after 1 second reaches its vertex at 20 feet. What is an equation that defines this function? Write the quadratic equation in vertex form and in the form $y=ax^2+bx+c$.

\begin{minipage}[t]{0.56\textwidth}
\vspace{0pt}
\begin{enumerate}[label=\textbf{\arabic*.}, itemsep=3pt]
\item The \textbf{vertex} is $(h,k)=(1,20)$, so\\ $y=a(x-1)^2+20$.
\item The initial height is the $y$-intercept, so $(x,y)=(0,4)$ is on the graph. Substitute it and solve for $a$:\\
$4=a(0-1)^2+20$, \ so \ $4=a+20$, \ $a=-16$
\item Vertex form: $\boxed{y=-16(x-1)^2+20}$
\item Expand to get $y=ax^2+bx+c$:\\
$y=-16(x^2-2x+1)+20=-16x^2+32x-16+20$\\
$\boxed{y=-16x^2+32x+4}$
\end{enumerate}
\textbf{Check:} $c=4$ is the initial height, and $a=-16<0$, so the parabola opens downward: the vertex is a maximum.
\end{minipage}\hfill
\begin{minipage}[t]{0.41\textwidth}
\vspace{0pt}
\begin{center}
\begin{tikzpicture}
\begin{axis}[sheetgrid, xmin=0, xmax=2.3, ymin=0, ymax=22.5,
  xtick={0,0.5,...,2}, ytick={0,4,...,20}, minor x tick num=1, minor y tick num=1,
  minor grid style={gray!25},
  xticklabels={0,,1,,2}, x=2.5cm, y=0.24cm,
  xlabel={Time (sec)}, ylabel={Height (ft)},
  xlabel style={at={(axis description cs:0.5,-0.1)}, anchor=north, font=\scriptsize},
  ylabel style={at={(axis description cs:-0.16,0.5)}, anchor=south, rotate=90, font=\scriptsize}]
\addplot[very thick, red, domain=0:2.1, samples=80, ->] {-16*(x-1)^2+20};
\node[font=\scriptsize, red, fill=white, inner sep=1pt] at (axis cs:1,9) {$y=ax^2+bx+c$};
\node[closed, red] at (axis cs:0,4) {};
\node[closed, red] at (axis cs:1,20) {};
\node[font=\scriptsize, fill=white, inner sep=1pt, anchor=north west] at (axis cs:0.06,3.4) {$(0,4)$ initial height};
\node[font=\scriptsize, fill=white, inner sep=1pt, anchor=south] at (axis cs:1,20.6) {vertex $(1,20)$};
\end{axis}
\end{tikzpicture}
\end{center}
\end{minipage}

\textbf{Say it:} The vertex is \blank[1.6cm], so $h=\blank[0.9cm]$ and $k=\blank[0.9cm]$. I substituted the point \blank[1.6cm] for $x$ and $y$ and got $a=\blank[1cm]$. The parabola opens \blank[2cm] because $a$ is \blank[2cm], so the vertex is a \blank[2.2cm].

\wspart{Try It 4}{Write the equation from the graph}
The graph shows the height of a ball with respect to time. What is the equation of the function? Write the equation in vertex form. Then write the equation in the form $y=ax^2+bx+c$.

\begin{minipage}[t]{0.44\textwidth}
\vspace{0pt}
\begin{tikzpicture}
\begin{axis}[sheetgrid, xmin=0, xmax=5.25, ymin=0, ymax=11.5,
  xtick={0,1,...,5}, ytick={0,2,...,10}, minor x tick num=1, minor y tick num=1,
  minor grid style={gray!45}, x=1.25cm, y=0.42cm,
  xlabel={Time (s)}, ylabel={Height (ft)},
  xlabel style={at={(axis description cs:0.5,-0.1)}, anchor=north, font=\scriptsize},
  ylabel style={at={(axis description cs:-0.12,0.5)}, anchor=south, rotate=90, font=\scriptsize}]
\addplot[very thick, red, domain=0:4.58, samples=80, ->] {-1.5*(x-2)^2+10};
\node[closed, red] at (axis cs:0,4) {};
\node[closed, red] at (axis cs:2,10) {};
\node[font=\scriptsize, fill=white, inner sep=1pt, anchor=north west] at (axis cs:0.1,3.6) {$(0,4)$};
\node[font=\scriptsize, fill=white, inner sep=1pt, anchor=south] at (axis cs:2,10.35) {$(2,10)$};
\end{axis}
\end{tikzpicture}
\end{minipage}\hfill
\begin{minipage}[t]{0.53\textwidth}
\vspace{0pt}
Vertex $(h,k)=\blank[2.2cm]$

Point $(x,y)=\blank[2.2cm]$

$y=a(x-\blank[0.8cm])^2+\blank[0.8cm]$

Substitute the point and solve for $a$:
\lines[3]
Vertex form: $y=\blank[4cm]$
\end{minipage}

Expand to get the form $y=ax^2+bx+c$:
\lines[2]
$y=\blank[5cm]$ \qquad Which number in this form is the initial height? \blank[2.5cm]

\newpage

\wspart{IXL}{Practice for the quiz}
Do the starred skills first. \textbf{SmartScore 70 or above on each starred ($\bigstar$) skill replaces your quiz score.}

\vspace{2pt}
\begin{tabularx}{\textwidth}{l Y Y}
\toprule
\textbf{Code} & \textbf{Skill} & \textbf{Feeds} \\
\midrule
\ixlrow{L8C $\bigstar$}{Characteristics of quadratic functions: equations}{Quiz 2: vertex, maximum or minimum from $a(x-h)^2+k$ (Example 2)}
\ixlrow{URV $\bigstar$}{Write a quadratic function in vertex form from its vertex and another point}{Quiz 3 and 4: vertex + one point gives $a$ (Example 4)}
\ixlrow{KQL}{Transformations of quadratic functions}{Optional: Quiz 1 and 5 (full Savvas quiz only)}
\ixlrow{WMS}{Characteristics of quadratic functions: graphs}{Optional warm-up: key features read off a drawn parabola}
\bottomrule
\end{tabularx}

\wspart{Exit question}{Amaya's serve (Practice 15)}
\fbox{\parbox{\dimexpr\textwidth-2\fboxsep-2\fboxrule}{%
Amaya is standing 30 ft from a volleyball net. The net is 8 ft high. Amaya serves the ball. The path of the ball is modeled by the equation $y=-0.02(x-18)^2+12$, where $x$ is the ball's horizontal distance in feet from Amaya's position and $y$ is the distance in feet from the ground to the ball.

\textbf{a.} How far away horizontally is the ball from Amaya when it is at its maximum height? Explain.

The vertex is \blank[2cm]. The parabola opens \blank[2cm] because $a=\blank[1.4cm]$, so the vertex is the \blank[2.2cm]. The ball is \blank[1.4cm] ft from Amaya because \lines[2]

\textbf{b.} Describe how you would find the ball's height when it crosses the net at $x=30$.
\lines[4]
}}

\vspace{6pt}
\textbf{Before the quiz, finish these sentences.}

In $y=a(x-h)^2+k$, the vertex is \blank[2cm] and the axis of symmetry is the line \blank[2cm].

When $a$ is negative, the parabola opens \blank[2.4cm] and $k$ is the \blank[2.4cm] value.

To find $a$, I substitute the \blank[2.4cm] for $h$ and $k$ and one other \blank[2.4cm] for $x$ and $y$.

\end{document}
